% D04-009. AI-04 P0060--P0070, P0180--P0254.
% Числа сверены с compare_4models_latest_20260610_204942.
\section{Сопоставление моделей при распознавании замыканий на землю}
\label{sec:ch4_ozz_validation}
\conceptsource{Статья 4. Гибридное обучение, RTDS и реальные подтверждённые/сложные отрицательные записи. Обобщённое сравнение и подробные примеры дополняют друг друга.}


\iffalse
(Примечание автора: Бегло просмотрел и даже не стал вычитывать. В исходных работах явно намного больше и понятнее, особенно больше картинок как раз, а у тебя прост слова и всё, без какого-то отображения практически. Я за то, чтобы максимально полно переносить статьи, пусть и перекомбинируя их и объединяя мысли) **[AA099: рассмотрено 27.10; ожидает проверки]**
(Примечание автора: Ок, закрываем примечание, если что будет – то уже локально.)
\fi
\authorreviewnote{AA099}

Задача распознавания ОЗЗ позволяет проверить общие
архитектурные решения на конкретном классе повреждений.
В ней доступны как моделируемые процессы с известным
способом формирования дуги, так и реальные записи
с подтверждённым наличием замыкания.
Их совместное использование позволяет разделить
способность воспроизводить модельные закономерности
и результат переноса на зарегистрированный материал.

\subsection{Подготовка опыта и единицы оценки}
\label{sec:ch4_ozz_test_protocol}
\conceptsource{Статья 4. Четыре модели дугового процесса, гибридные данные, оконные и файловые показатели.}


В исследовании трансформерных моделей для детекции событий\footnote{Евдаков А.Е., Яблоков А.А., Филатова Г.А.
Практическое применение физически-информированных нейронных
сетей для распознавания сложных аварийных режимов электроустановок.}
использована гибридная обучающая выборка.
Её основу составили процессы для сети 10~кВ
с изолированной нейтралью, полученные на RTDS.
Варьировались параметры цепи и условия возникновения
повреждения. Рассматривались устойчивое ОЗЗ
и три варианта перемежающегося дугового процесса,
заданные моделями Петерсена, Петерса--Слепяна и Белякова.
Реальные осциллограммы без ОЗЗ использовались
как примеры переходного фона с отсутствующими
целевыми признаками этих повреждений.

Для спектральных вариантов анализируется окно
длительностью десять периодов промышленной частоты.
Окно Фурье и временной шаг задаются относительно периода,
что позволяет работать с разными частотами дискретизации
исходных записей. При частоте 50~Гц шаг в одну восьмую
периода соответствует 2,5~мс.
Входные величины и физические преобразования
рассмотрены в разделах~\ref{sec:ch4_signal_representations}
и~\ref{sec:ch4_physkan_architectures}.

Валидационная часть RTDS содержит 960 файлов,
по 240 для каждого типа, и 116\,160 окон анализа.
Для неё оцениваются четыре двоичных выхода.
Macro-F1 усредняет их F1, а точное совпадение
означает одновременное совпадение всех четырёх меток
в проверяемой единице. При оконной оценке для каждого
класса берётся максимум выходной оценки по зонам.
Целевая метка окна также определяется максимумом
целевых меток его зон. После сравнения с порогом
0,5 это отвечает на вопрос, присутствует ли событие
хотя бы в одной зоне окна. Правильность его положения
внутри окна при таком подсчёте отдельно не оценивается.

Для позонной оценки сравниваются метки каждой выходной
зоны каждого окна. При перекрывающихся окнах один
физический участок может участвовать в подсчёте
несколько раз. Эти показатели не являются оценкой
полностью независимых интервалов или долей
распознанных файлов.

% D04-011 R4-11 BEGIN
\ifdefined\AuthorReview
\par\smallskip\begingroup\small\raggedright\color{red!60!black}
\noindent\textbf{Примечание рецензента [R4-11], существенно.} Явно назвать правило оконной оценки: для каждого класса берётся максимум по зонам, то есть «событие было хотя бы где-нибудь в окне». Это не качество позонной локализации. В сохранённом отчёте Physical KAN оконный Macro-F1=0,978947, а позонный=0,946610. Позонные показатели уже есть для всех четырёх моделей, новое обучение не требуется. Их стоит кратко показать рядом, если вывод главы относится к точному выделению интервалов.\par\endgroup\smallskip
\fi
% D04-011 R4-11 END
% D04-012-A RESPONSE R4-11 BEGIN
\ifdefined\AuthorReview
\par\smallskip\begingroup\small\raggedright\color{blue!55!black}
\noindent\textbf{Ответ на R4-11, D04-012-A.}
Описаны максимум по зонам для оценок и целей,
порог 0,5 и повторный учёт участков перекрывающихся окон.
В таблице~\ref{tab:ch4_ozz_units} рядом приведены
оконные и позонные Macro-F1 из тех же четырёх отчётов.
Вывод о наибольшем результате разделён по единицам оценки.
Исправление внесено, ожидает повторной проверки.
\par\endgroup\smallskip
\fi
% D04-012-A RESPONSE R4-11 END

(Примечание автора: Хмм… Проверит рецензент и закроет примечание.)

На реальном материале отдельно рассматриваются
829 пригодных записей с подтверждённым ОЗЗ
и 1004 пригодные сложные записи без ОЗЗ.
Во вторую группу входят случаи, ранее вызвавшие
ложное обнаружение базовым алгоритмом.
Это специально отобранный материал для проверки
различения похожих переходных процессов.
Его состав не воспроизводит доли режимов
в непрерывном эксплуатационном потоке.

На файловом уровне наличие ОЗЗ определяется
хотя бы одним положительным решением соответствующих
выходов за время записи. Для подтверждённых ОЗЗ
вычисляется доля обнаружений, для второй группы
доля ложных обнаружений.
Один краткий лишний отклик уже может сделать
отрицательный файл положительным по такому правилу.
Поэтому файловые доли рассматриваются вместе
с временной разметкой, а не вместо неё.

\subsection{Позонное распознавание моделируемых процессов}
\label{sec:ch4_rtds_examples}
\conceptsource{Статья 4, рисунок 2. Исходные примеры четырёх
моделируемых процессов показывают связь сигналов, целевого
интервала и выходов модели. Они дополняют сводные показатели.}

(Примечание автора: Подожди, может я что-то путаю, а не повторяется ли с позонной оценкой чуть раньше? Вроде бы совсем недавно было похожее, прост на других моделях. Хотя и текста тут почти нет, не совсем понял, что и зачем так в главу выделять?)
В моделируемой записи известны как способ формирования
повреждения, так и заданный временной интервал.
На рисунках~\ref{fig:ch4_rtds_examples_a}--\ref{fig:ch4_rtds_examples_b}
показаны четыре характерных процесса.
Верхние поля содержат токи и напряжения,
ниже расположены оценки четырёх выходов ОЗЗ.
Закрашенный целевой интервал задаётся при формировании
примера и позволяет сопоставить ответ модели
с наличием повреждения во времени.

Для устойчивого повреждения сохраняется длительный
положительный выход соответствующего класса.
При перемежающемся процессе изменение оценок
сопоставляется с повторяющимися нарушениями формы сигнала.
Выход для всего окна скрыл бы различия начала,
продолжения и завершения этих участков.
Позонное представление позволяет их рассматривать отдельно.

При этом высокий отклик на выбранном синтетическом примере
характеризует согласие с заданной моделью процесса.
Следующий этап включает другие комбинации параметров
и реальные записи. Именно совместная проверка показывает,
какие признаки сохранили полезность за пределами
идеализированного сигнала.

% Четыре процесса выводятся вместе перед численным сравнением.
\begin{figure}[!htbp]
 \centering
 \includegraphics[width=\textwidth,height=0.79\textheight,keepaspectratio]{ch04_rtds_stable_article.png}
 \caption{Позонные выходы Physical KAN-Transformer
 на моделируемом устойчивом замыкании на землю}
 \label{fig:ch4_rtds_examples_a}
\end{figure}
\begin{figure}[p]
 \centering
 \includegraphics[width=\textwidth,height=0.79\textheight,keepaspectratio]{ch04_rtds_petersen_article.png}
 \caption{Позонные выходы на процессе
 по модели Петерсена}
 \label{fig:ch4_rtds_petersen}
\end{figure}
\begin{figure}[p]
 \centering
 \includegraphics[width=\textwidth,height=0.79\textheight,keepaspectratio]{ch04_rtds_peters_slepyan_article.png}
 \caption{Позонные выходы на процессе
 по модели Петерса--Слепяна}
 \label{fig:ch4_rtds_peters_slepyan}
\end{figure}
\begin{figure}[p]
 \centering
 \includegraphics[width=\textwidth,height=0.79\textheight,keepaspectratio]{ch04_rtds_belyakov_article.png}
 \caption{Позонные выходы на процессе
 по модели Белякова}
 \label{fig:ch4_rtds_examples_b}
\end{figure}
\figurereview{Четыре исходных поля рисунка 2 статьи 4. При следующей
редакции увеличить подписи токов, напряжений и выходов, перевести
Stable/Petersen/PetersSlepyan/Belyakov. Сейчас включены исходные
изображения без перерисовки. Это иллюстрации отдельных процессов,
они не заменяют обновлённую сравнительную таблицу.}

\clearpage
\subsection{Результаты моделируемых и реальных процессов}
\label{sec:ch4_ozz_comparison}
\conceptsource{Статья 4, обновлённая сравнительная серия. Совместно приведены RTDS, чувствительность к реальным ОЗЗ и ложные обнаружения на сложных отрицательных записях.}


В таблице~\ref{tab:ch4_ozz_comparison} приведено
сопоставление четырёх конфигураций после 100 эпох
целевого обучения распознаванию ОЗЗ.
Для каждой конфигурации оценивается состояние
после завершения обучения. Условия проверки
и единицы расчёта показателей сохраняются одинаковыми
в пределах сравниваемой серии.

(Примечание автора: Ок, вроде бы всё исправлено, закрывай.)

% D04-011 R4-10 BEGIN
\ifdefined\AuthorReview
\par\smallskip\begingroup\small\raggedright\color{red!60!black}
\noindent\textbf{Примечание рецензента [R4-10], исправить до приёмки.} Подтверждённое расхождение: в сохранённых sim\_ozz\_evaluation.json, соответствующих четырём строкам этой таблицы, config.epochs равен 100. В каталогах запусков summary.json также содержит total\_epochs=100, а журналы заканчиваются epoch=99. Числа таблицы совпадают именно с этими отчётами. Заявление о 200 эпохах из статьи нельзя оставлять без отдельного подтверждения другой серии. У физических моделей также расходится число параметров в сводке обучения и отчёте оценки: требуется зафиксировать историческую версию, не загружая автоматически нынешний код как прежнюю модель.\par\endgroup\smallskip
\fi
% D04-011 R4-10 END
% D04-012-A RESPONSE R4-10 BEGIN
\ifdefined\AuthorReview
\par\smallskip\begingroup\small\raggedright\color{blue!55!black}
\noindent\textbf{Ответ на R4-10, D04-012-A.}
Число эпох исправлено на 100 по четырём отчётам оценки,
конфигурациям запусков, сводкам и журналам с последней
эпохой 99 при нумерации с нуля. Прежнее утверждение
о 200 эпохах было ошибочным. Расхождение числа параметров
физических моделей между сводками обучения и оценкой
остаётся открытым вопросом исторической реализации.
Код и веса не менялись, новое применение моделей не выполнялось.
\par\endgroup\smallskip
\fi
% D04-012-A RESPONSE R4-10 END

\begin{table}[htbp]
 \centering
 \small
 \caption{Оконная оценка на RTDS и обнаружение ОЗЗ
 в реальных записях}
 \label{tab:ch4_ozz_comparison}
 \begin{tabular}{p{4.0cm}rrrr}
  \hline
  Модель & \multicolumn{2}{c}{RTDS} &
  \multicolumn{2}{c}{Реальные записи, \%} \\
   & Macro-F1 & Совпадение & Обнаружено & Ложно \\
  \hline
  Physical KAN & 0,9789 & 0,9719 & 94,81 & 63,65 \\
  Physical MLP & 0,9894 & 0,9833 & 93,37 & 67,43 \\
  Базовая спектральная & 0,9923 & 0,9864 & 92,64 & 55,68 \\
  Базовая на отсчётах & 0,9914 & 0,9857 & 95,30 & 63,15 \\
  \hline
 \end{tabular}
\end{table}

При оконной оценке RTDS наибольший Macro-F1 имеет базовая спектральная
модель. В конфигурации Physical KAN он несколько ниже,
при этом значения F1 отдельных классов различаются.
Для устойчивого ОЗЗ получено 0,9939,
для моделей Петерсена и Петерса--Слепяна 0,9802 и 0,9831,
для модели Белякова 0,9586.
В этой серии включение физических блоков
не даёт наибольшей средней оконной оценки на моделируемых данных.

В таблице~\ref{tab:ch4_ozz_units} сопоставлены
обе единицы оценки при пороге 0,5.
Позонный Macro-F1 ниже оконного у всех четырёх моделей.
Наибольшее позонное значение получено у Physical MLP,
а не у базовой спектральной модели.
Поэтому результат обнаружения события в окне
рассматривается отдельно от правильности его разметки
по выходным зонам.

\begin{table}[htbp]
 \centering\small
 \caption{Macro-F1 при оконной и позонной оценке
 тех же конфигураций на RTDS}
 \label{tab:ch4_ozz_units}
 \begin{tabular}{lrr}
  \toprule
  Модель & По окнам & По зонам \\
  \midrule
  Physical KAN & 0,9789 & 0,9466 \\
  Physical MLP & 0,9894 & 0,9647 \\
  Базовая спектральная & 0,9923 & 0,9532 \\
  Базовая на отсчётах & 0,9914 & 0,9494 \\
  \bottomrule
 \end{tabular}
\end{table}

На реальных подтверждённых ОЗЗ Physical KAN обнаруживает
786 из 829 записей, пропуская 43.
На сложных записях без ОЗЗ положительный выход
появляется в 639 из 1004 файлов.
Соответственно, 365 записей этой группы
не вызывают обнаружения.
Большая чувствительность здесь сочетается
с существенной ложной активностью.

При сравнении четырёх конфигураций наибольшую
чувствительность имеет модель на отсчётах,
а наименьшую долю ложных обнаружений
базовая спектральная модель.
Physical KAN даёт меньшую ложную активность,
чем Physical MLP, но не меньшую, чем все варианты.
По этим результатам нельзя выделить одну конфигурацию,
одновременно превосходящую остальные по обоим критериям.

В реальной проверке для всех четырёх выходов
используется порог 0,5.
Отдельная настройка порога по валидации
рассматривается отдельно и не приписывается
этим файловым долям.
Выбор порогов по проверяемому реальному архиву
также изменил бы назначение этого архива
в процедуре оценки.

\subsection{Характер ошибок и трактовка физических классов}
\label{sec:ch4_ozz_errors}
\conceptsource{Статья 4. Содержание ложных выходов и границы физической трактовки классификации дуги.}


В Physical KAN наиболее частый выход на сложных
отрицательных записях соответствует модели Белякова.
Он присутствует в 519 из 1004 файлов.
Для устойчивого ОЗЗ отмечены 140 файлов,
для Петерсена 31, для Петерса--Слепяна 170.
Эти группы пересекаются, поэтому их числа
не складываются в число ошибочных записей.

Такая картина указывает, какую часть распознавания
следует дополнительно проверять на коммутационных
и иных неаварийных переходных процессах.
Она не устанавливает единственную физическую причину
ложного выхода. В частности, сходство кратких искажений
с моделируемым дуговым процессом требует
предметного разбора выбранных осциллограмм.

На реальных подтверждённых ОЗЗ несколько выходов
также могут появляться в одной записи.
Однако подтверждённое наличие ОЗЗ и проверенная
разметка конкретной теории горения дуги различаются.
В используемом реальном архиве последняя
не задана как независимый эталон.
Поэтому распределение выходов описывает решения модели,
а не доказанную распространённость физических механизмов.
Последовательность разных выходов полезна
для отбора примеров последующего анализа.

\subsection{Изменение выходов на реальных осциллограммах}
\label{sec:ch4_ozz_real_examples}
\conceptsource{Статья 4, исходные примеры рисунка 6. Показаны
меняющиеся выходы при затухающем процессе и длительный отклик
на другой подтверждённой записи. Файловое подтверждение ОЗЗ
не превращается в эталонную разметку каждой зоны и теории дуги.}

На реальной записи характер процесса может изменяться
на протяжении одного события.
В примере на рисунке~\ref{fig:ch4_real_ozz_fading}
вместе с изменением токов и напряжений меняются
оценки нескольких классов. Одна протяжённая файловая
метка ОЗЗ не описывает эти изменения.
Позонные выходы позволяют выбрать для дальнейшего
разбора начало процесса, повторные изменения
и его затухающий участок.

\begin{figure}[p]
 \centering
 \includegraphics[width=\textwidth]{ch04_real_ozz_fading_article.png}
 \caption{Позонные выходы на реальной записи
 с изменяющимся и затухающим процессом ОЗЗ}
 \label{fig:ch4_real_ozz_fading}
\end{figure}
\figurereview{Исходный пример 6в статьи 4. Перевести подписи
четырёх выходов. Последовательность их активации показывается
как ответ сети; без независимой разметки не подписывать её
как доказанную смену физических теорий горения дуги.}

Другой пример приведён на
рисунке~\ref{fig:ch4_real_ozz_long_response}.
Он полезен для рассмотрения длительного отклика
одного из выходов за пределами наиболее заметного
начального переходного участка.
По общему положительному ответу на запись нельзя
определить, какая часть такого интервала верна.
Для этого рассматриваются сами измеренные сигналы
и уточняется временная разметка.

\begin{figure}[p]
 \centering
 \includegraphics[width=\textwidth]{ch04_real_switching_error_article.png}
 \caption{Пример длительного выхода модели
 на отдельной реальной записи с подтверждённым ОЗЗ}
 \label{fig:ch4_real_ozz_long_response}
\end{figure}
\figurereview{Исходный пример 6д статьи 4. Автор отмечал, что часть
длительного выхода Белякова выглядит ошибочной. Сохранить это
для предметного разбора границ, не подписывать весь файл как
отрицательный: в исходнике он относится к подтверждённым ОЗЗ.}

Таким образом, количественная проверка по файлам
и подробное рассмотрение выходов решают разные задачи.
Первая показывает, обнаружено ли повреждение хотя бы раз.
Второе позволяет исследовать протяжённость и устойчивость
сигнала, реакцию на отдельные изменения режима
и участки, для которых требуется уточнение разметки.

% Сначала вывести разобранные реальные осциллограммы.
\clearpage
\subsection{Вычислительные затраты и значение для БАВР}
\label{sec:ch4_ozz_runtime}
\conceptsource{Статья 4 и связь с главой 5. Вычислительный бюджет и ограничения прямого использования детекции в БАВР.}

% D04-011 R4-12 BEGIN
\ifdefined\AuthorReview
\par\smallskip\begingroup\small\raggedright\color{red!60!black}
\noindent\textbf{Примечание рецензента [R4-12], уточнить.} Добавить условия измерения CPU для этой серии и явно отделить их от Intel Core Ultra 7 165U из исследования признаков. Для инженерного вывода существенны также частота выдачи решений и возможность обработки без накопления очереди. Шаг спектральных признаков 2,5 мс сам по себе не означает запуск всей сети каждые 2,5 мс. По 23,62 мс одного вызова нельзя ни подтвердить работу в таком темпе, ни вычислить полное время действия БАВР.\par\endgroup\smallskip
\fi
% D04-011 R4-12 END
\ifdefined\AuthorReview
\par\smallskip\begingroup\small\raggedright\color{blue!55!black}
\noindent\textbf{Ответ на R4-12, D04-012-B.} Добавлены подтверждённые условия функции измерения: один вход, 10 прогревов, среднее 100 повторов, без вычисления спектра. Процессор и число потоков для исторической серии не восстановлены, Intel Core Ultra 7 165U ей не приписан. Шаг признаков не назван частотой запуска сети; аппаратный вопрос остаётся открытым.\par\endgroup\smallskip
\fi

(Примечание автора: Хм…. Ок, пусть будет так, только в рамках всей диссертации надо как-то быть аккуратным с тем, что запуски получается были на разных устройствах, моём домашнем, и некотором внешнем. Вот этот момент не упусти и правильно следи за этим, можно в задачу куда-то даже, чтобы и в будущем не забыть.)

В той же серии измеренное время расчёта окна на CPU
составляет 23,62~мс для Physical KAN,
12,55~мс для Physical MLP,
5,34~мс для базовой спектральной модели
и 5,68~мс для модели на отсчётах.
Измеряется среднее время одного вызова сети
после десяти прогревочных вызовов по ста повторениям.
На вход уже подан сформированный массив,
предварительное спектральное преобразование в это время не входит.
Полная аппаратная конфигурация этой серии не зафиксирована
в использованных сводках. Поэтому значения сопоставляются
внутри серии и не приписываются процессору из исследования
признакового пространства.
Они не являются временем полного действия автоматики.
Дополнительно учитываются накопление сигнала,
подготовка признаков, выдача решения
и аппаратная реализация.

Сравнение этих конфигураций характеризует модели целиком.
Физический MLP-вариант меняет несколько частей обработки,
поэтому разницу нельзя целиком приписать
одной замене функций KAN на MLP.
Для выделения вклада отдельного узла используется
сопоставимый опыт с фиксированными остальными частями,
как пояснено в разделе~\ref{sec:ch4_physical_validation}.

Для БАВР распознавание ОЗЗ полезно как часть описания
состояния сети и последовательности переходного процесса.
Результаты показывают также необходимость проверки
на похожих режимах без замыкания.
Высокая оценка на моделируемых процессах и обнаружение
большинства реальных ОЗЗ сами по себе не обосновывают
непосредственное использование выхода как разрешения
или блокировки переключения.
Связь с направлением мощности и логикой действия
рассматривается в специальной задаче следующей главы.

